\section{问题四：双引擎情景演算、右删失处理与有限库存机制}
\noindent\textbf{本问结论：}链A混合入沟体积$2.635835\times10^6$ m$^3$；D8基线32.766603 min到达口岸代理，16个配置全部在26045.494670 m分析域末端右删失；一维基线在9.336769 km自然停止，未到口岸。质量相对残差$2.075820\times10^{-15}$只说明离散记账闭合，不构成地形、流变或现场过程验证。

\subsection{任务目标与模型定位}
问题四要求在给定灾前地形、源区体积与传播参数下，估计过程沿沟演进、口岸到达时间、运程和影响强度。现有资料缺少真实释放过程、沿程侵蚀物源、实测流深流速和现场流变，故本文建立参数透明的条件情景模型，而不把结果称为已率定的灾害重建。两个引擎承担不同任务：一维纵剖面用于检验准恒定断面下的自然停止与沿程尺度；D8网络用于显式连续性推进，输出首达场、历史峰值、终态场和质量收支。D8每格只向一个下游格路由，不含横向动量、漫滩扩散和完整浅水方程\cite{obrien1993}。

两个引擎使用同一公共注入断面和同一口岸代理断面，并以注入断面统一计时。口岸代理是题面地点在DEM上的吸附断面，不是实测受灾断面；全域出口是另一对象，不能代替口岸断面。两引擎结果差异用于揭示结构不确定性，不是两次独立观测，也不因D8到达口岸而删除一维自然停止的结果。

\subsection{公共几何、源区体积与注入边界}
在附件4原生冰川网格上，以冰厚至少50 m且坡度至少25$^\circ$筛选源区，共335个原生像元。先在原生网格积分，再映射到演算网格；335次投影全部落域、去重为323格，体积不在映射阶段重复计数。链A的体积链为
\begin{equation}
V_{\mathrm{src}}=\sum_{k}h_{\mathrm{ice},k}a_k,\qquad
V_{\mathrm{rel}}=\kappa_{\mathrm{rel}}V_{\mathrm{src}},
\end{equation}
\begin{equation}
V_{\mathrm{df}}^{(A)}
=\gamma\frac{\rho_{\mathrm{ice}}}{\rho_{\mathrm{df}}}V_{\mathrm{rel}},
\qquad
\rho_{\mathrm{df}}
=C_v\rho_s+(1-C_v)\rho_w.
\end{equation}
取$\kappa_{\mathrm{rel}}=0.10$、$\gamma=3$、$\rho_{\mathrm{ice}}=900$ kg/m$^3$、$\rho_s=2650$ kg/m$^3$、$\rho_w=1000$ kg/m$^3$、$C_v=0.45$，则$\rho_{\mathrm{df}}=1742.5$ kg/m$^3$。源冰总量为$1.70109\times10^7$ m$^3$，释放冰量为$1.70109\times10^6$ m$^3$，链A混合入沟体积为
\begin{equation}
V_{\mathrm{target}}^{(A)}
=2635835.294118\ \mathrm{m^3}.
\end{equation}

链B以厚冰筛选区的实际平均厚度、面积和可动比例形成$1.123497\times10^7$ m$^3$的大体量对照。它不经过与链A相同的密度与物源换算，只表示本情景族的高体积端点，不是现实严格上界。比较两链前需统一对象是源冰、释放冰还是混合入沟体积，避免把不同对象混为同一规模。

链A落域格中汇流累积最大者作为公共注入点，口岸代理取该下行路线距口岸工作坐标最近的有效格，吸附偏移5.186582 m。中心线水平长度20.882237 km、总落差2126.790527 m，重采样为209个约100 m节点。原始纵剖面含137个上升步，可能来自真实局部地形、填洼或重采样；本文仅以单调不增回归构造演算剖面，最大调整量39.517428 m。该处理便于Manning局部坡度计算，也会抹平部分蓄滞，不应视为仅修正地形。

\subsection{一维诊断引擎与自然停止}
以特征流量$Q_{\mathrm{char}}=V_0/T_{\mathrm{pass}}$驱动准恒定Manning断面\cite{chow1959}：
\begin{equation}
B(s)=B_0(1+0.8s/L).
\end{equation}
\begin{equation}
h(s)=\left[\frac{nQ_{\mathrm{char}}}{B(s)\sqrt{S_{\mathrm{eff}}(s)}}\right]^{3/5},
\qquad
v(s)=\frac{Q_{\mathrm{char}}}{B(s)h(s)}.
\end{equation}
基线取$V_0=2635835.294118$ m$^3$、$B_0=256.880875$ m、$n=0.10$ s/m$^{1/3}$、$T_{\mathrm{pass}}=900$ s、$\mu_H=0.145259$、$S_{\mathrm{eff}}\ge0.01$。经验能量线为
\begin{equation}
E(s)=E_0+2g\left[(z_0-z(s))-\mu_H(s-s_0)\right],
\qquad E_0=200\ \mathrm{m^2/s^2}.
\end{equation}
连续5个节点满足$E\le0$时判自然停止。$E$具有速度平方量纲，是准稳断面可达性的诊断量，不是完整瞬态能量守恒解；停止规则本身也属于结构假设。

基线在9.336769 km处自然停止，停止时刻1783.858032 s（29.73 min），未达20.882237 km口岸，状态为自然能量停止。17个去重一维配置中16个未达口岸，仅$\mu_H=0.1017$的条件配置于83.3575 min触及口岸末端；后者是口岸右截断下界，不是自然停止距离。该负结果原样保留，不通过删除未达配置或反向挑选摩擦系数制造统一到达。基线沿程剖面见图\ref{fig:q4-one-d}；能量线在上游判停，图中线段不延伸为虚假的口岸到达时间。

\paperfig{figures/04_沿程演进.png}{0.94}{一维沿程演进。灰色区为自然停止后的未到达段，橙色虚线为口岸代理}{fig:q4-one-d}

强度采用
\begin{equation}
I=\rho_{\mathrm{df}}hv^2\quad[\mathrm{N/m}],
\end{equation}
其中$\rho_{\mathrm{df}}=1742.5$ kg/m$^3$。该量表示单位宽度动量通量尺度，不是压强。基线沿程最大强度为153.092248517 kN/m，只作情景分级，不作为本地已验证的建筑破坏阈值。基线未到口岸，故不在口岸生成物理解释上的堆积扇体；任何几何扇形只能作为额外条件包络。

\subsection{D8屈服门槛—Manning路由与临界量}
在Priority-Flood填洼地形上，每格按米制最大下坡选择唯一接收者\cite{tarboton1997}。记格面积$a_i$、等效过流宽度$W_i$和源项$q_i^{\mathrm{src}}$，连续方程与代数流量为
\begin{equation}
\frac{dh_i}{dt}
=\frac{\sum_{j:r(j)=i}Q_j-Q_i}{a_i}+q_i^{\mathrm{src}},
\end{equation}
\begin{equation}
Q_i=
\begin{cases}
n^{-1}h_i^{5/3}\sqrt{S_{\mathrm{eff},i}}W_i,
&\rho_{\mathrm{df}}gh_iS_{\mathrm{eff},i}\ge\tau_y,\\
0,&\text{其他},
\end{cases}
\qquad
S_{\mathrm{eff},i}=\max(S_i,S_{\min}).
\end{equation}
这是屈服开关加Manning运动波，不含独立库仑摩擦和宾汉塑性黏度，不是完整的库仑—宾汉流变。

由屈服条件可写
\begin{equation}
h_{y,i}S_{\mathrm{eff},i}
=\frac{\tau_y}{\rho_{\mathrm{df}}g}.
\end{equation}
右端在固定屈服门槛与密度下具有长度量纲，但它是“屈服水深×坡度”的乘积，不是可直接解释为水深的$h_{y}$。记该乘积为$\Pi_c$，基线值为
\begin{equation}
\Pi_c=h_{y}S_{\mathrm{eff}}
=0.04680025213635838\ \mathrm{m}.
\end{equation}
坡度无量纲，故单位仍为m。按口岸代理的有效坡度$S_{\mathrm{eff}}=0.02$换算，得到口岸坡度的等效屈服水深
\begin{equation}
h_{y,\mathrm{port}}
=\frac{\Pi_c}{0.02}
=2.3400126068179192\ \mathrm{m}.
\end{equation}
前者只对应局部临界水深，后者只对应口岸坡度，均不作为全线统一门槛。

网络基准参数见表\ref{tab:q4-routing}。参数为给定情景值，不是独立率定后验。

\begin{table}[htbp]\centering\small
\caption{D8网络路由基准参数}\label{tab:q4-routing}
\begin{tabular}{@{}l c c l@{}}
\toprule
参数&基准值&单位&含义\\
\midrule
$V_{\mathrm{target}}$&2635835.294118&m$^3$&链A混合入沟体积\\
$n$&0.10&s/m$^{1/3}$&等效Manning糙率\\
$\tau_y$&800&Pa&屈服门槛\\
$C_v$&0.45&1&岩相体积分数情景\\
$S_{\min}$&0.02&1&缓坡驱动下限\\
$\tau_{\mathrm{src}}$&1800&s&三角源项历时\\
$h_{\min}$&0.10&m&首次到达分析阈值\\
$C_{\mathrm{CFL}}$&0.2&1&设定Courant上限\\
$T_{\mathrm{end}}$&21600&s&最长演算时限\\
\bottomrule
\end{tabular}
\end{table}

源项取
\begin{equation}
Q_{\mathrm{src}}(t)=
\begin{cases}
\dfrac{2V_{\mathrm{target}}}{\tau_{\mathrm{src}}}
\left(1-\dfrac{t}{\tau_{\mathrm{src}}}\right),&0\le t\le\tau_{\mathrm{src}},\\
0,&\text{其他},
\end{cases}
\end{equation}
并在每步直接使用累计函数
\begin{equation}
F(t)=V_{\mathrm{target}}(2u-u^2),
\qquad
u=\min\{\max(t/\tau_{\mathrm{src}},0),1\}
\end{equation}
之差注入，避免跨越源项结束点时多计体积。

同一通量从供体扣除、向受体加入，域外通量单独累计。固定宽度下特征波速$c=5v/3$，实际Courant数及步长约束为
\begin{equation}
Co_i=\frac{c_iW_i\Delta t}{a_i}\le C_{\mathrm{CFL}},
\qquad
\Delta t\le
C_{\mathrm{CFL}}\min_i\frac{a_i}{c_iW_i}.
\end{equation}
候选步先计入解析源量并估算源后波速，若超限则缩步重算。正性限制与CFL通过都不替代空间分辨率和物理参数验证。

\subsection{D8过程、峰值与域末端右删失}
首次到达定义为
\begin{equation}
t_{\mathrm{arr},i}=\inf\{t:h_i(t)\ge h_{\min}\},
\end{equation}
历史峰值逐步独立累积为
\begin{equation}
h_{\mathrm{peak},i}=\max_t h_i(t),\quad
v_{\mathrm{peak},i}=\max_t v_i(t),\quad
I_{\mathrm{peak},i}=\max_t\left[\rho_{\mathrm{df}}h_i(t)v_i(t)^2\right],
\end{equation}
峰值需由同一时刻的$h,v$计算，不用终态水深代替，也不拼接不同时刻的$\max h$与$\max v$。

基线口岸代理首达1965.996174 s，即32.7666029 min。16个去重D8配置全部到达口岸代理，首达范围为17.5704376--44.7278067 min；最快端由链B高体积情景给出，最慢端由$V_{\mathrm{low}}$体积情景给出。全部配置均在注入结束后转为静止，并因物质前锋到达或越过分析域末端而终止。两项需联合解释：注入后不再产生有效变化，不等于发生了物理自然停积；前锋抵达域边界，只说明演算被计算域截断。

全部配置的分析域末端里程均为26045.4946703 m。这只是计算域的右删失边界，不是自然最大运程、现场停积边界或危险区外缘，也不用于估计分析域之外的真实影响范围。

基线口岸历史峰深10.08530186 m、终态深度2.298002229 m，全域历史最大深13.31272084 m，口岸峰速6.601474228 m/s。终态口岸深度虽接近按口岸坡度换算的2.3400126068179192 m临界量，但低于该量并不表示永久稳定，未跨域演算也无法恢复真实停积位置。过程量和全部配置的终止条件见表\ref{tab:q4-results}，到达时间场与终态泥深场见图\ref{fig:q4-arrival}，过程快照见图\ref{fig:q4-snapshots}。快照只表示给定源项和参数下的条件演进，不含横向二维动力；终态仅为分析截止时的像元状态，图幅不表示自然最大运程或实测漫滩范围。

\paperfig{figures/04b_到达与停积.png}{0.92}{首达时间与终态泥深。红色星号为口岸代理}{fig:q4-arrival}

\paperfig{figures/04b_过程快照.png}{0.86}{D8过程快照（10、20、30、60 min）}{fig:q4-snapshots}

\begin{table}[htbp]\centering\small
\caption{D8基线结果}\label{tab:q4-results}
\begin{tabular}{@{}l l@{}}
\toprule
指标&当前结果\\
\midrule
口岸代理首达时间&1965.996174 s，即32.7666029 min\\
16配置首达范围&17.5704376--44.7278067 min\\
口岸历史峰深&10.08530186 m\\
口岸终态深度&2.298002229 m\\
全域历史最大深&13.31272084 m\\
口岸峰值速度&6.601474228 m/s\\
质量相对残差&2.075820191829099$\times10^{-15}$\\
最大实际Courant数&0.188281359\\
终止条件&注入结束后演算转为静止（全部配置）\\
停积性质&域末端右删失，非自然停积（全部配置）\\
域末端里程&26045.4946703 m\\
\bottomrule
\end{tabular}
\end{table}

\subsection{质量收支与数值闭合}
逐步记录累计注入、域内、出域和损失体积，质量账本为
\begin{equation}
V_{\mathrm{inj}}(t)-V_{\mathrm{dom}}(t)-V_{\mathrm{out}}(t)
=V_{\mathrm{loss}}(t).
\end{equation}
以实际注入体积为分母，基线最大相对残差为2.075820191829099$\times10^{-15}$；源项相对误差为0，全域最小深度非负，最大实际Courant数为0.188281359，未发生供体体积截断步骤。质量收支曲线见图\ref{fig:mass-ledger}。

这些数值说明源项积分、离散转移和状态记录在机器精度内闭合，CFL实现未超过给定上限；这些数值只说明数值实现一致；固定地形、单流向、屈服门槛、等效糙率和集中源项是否等于真实物理过程，仍需现场检验。

\subsection{有限库存、静止混合与分段动力对照}
一维和D8均以集中注入和固定流变处理物源；沿程静止物料加入后团块会加速还是减速，需要另设模型。为此增加固定路径代表团块模型，仅保留动摩擦$\mu_{\mathrm{dyn}}$、二次阻力$\kappa$和入口速度$v_0$三个全局工程量。它不替代前述引擎，其速度也不等于真实前锋速度或流深。

对固定水平路线，令相邻支持点三维段长为$L_k$、下坡角为$\alpha_k$，$w=v^2$，则
\begin{equation}
v\frac{dv}{d\ell}
=g(\sin\alpha_k-\mu_{\mathrm{dyn}}\cos\alpha_k)-\kappa v^2,
\qquad
\frac{dw}{d\ell}=2a_k-2\kappa w.
\end{equation}
左侧$v$因子不可省略。常参数段在$\kappa>0$时满足
\begin{equation}
w(x)=w_{\mathrm{in}}e^{-2\kappa x}
+\frac{a}{\kappa}(1-e^{-2\kappa x}),
\qquad
a=g(\sin\alpha-\mu_{\mathrm{dyn}}\cos\alpha).
\end{equation}
当$a<0$时首次停止距离为
\begin{equation}
x_{\mathrm{stop}}
=\frac{\log\left(1+\kappa w_{\mathrm{in}}/|a|\right)}{2\kappa},
\end{equation}
$\kappa=0$时取连续极限$x_{\mathrm{stop}}=w_{\mathrm{in}}/(2|a|)$。段内停止后不自动重启；未达目标是无后续输入、无再启动假设下的条件结论。

真实GIS未确认沿程可接触物源节点及库存，故真实路线最低基线为无新增节点。原始、填洼和单调三种高程剖面在固定基准$\mu_{\mathrm{dyn}}=0.10$、$\kappa=10^{-5}$ m$^{-1}$、$v_0=30$ m/s下，分别于18.360547、18.226423和18.535479 km首次停止，对应273.057496、281.660354和287.613654 s，均未到口岸。无节点工程扫描共144个配置，其中每个剖面48个配置；三种剖面合计48个配置正速通过、96个首次停止。该比例只描述已检验工程域，不是发生概率。

合成节点固定于路线水平里程的1/3和2/3处，仅作机制示范。节点库存及静止加入的动量、混合损失为
\begin{equation}
B_{r,j}=q_rV_{r,0}\pi_j,\qquad
B_{w,j}=q_wV_{w,0}\pi_j,\qquad
\sum_j\pi_j=1,
\end{equation}
\begin{equation}
m^+v^+=m^-v^-,\qquad
v^+=\frac{v^-}{1+\lambda},
\qquad
E_{\mathrm{mix}}
=\frac12m^-(v^-)^2\frac{\lambda}{1+\lambda}\ge0,
\end{equation}
其中$\lambda=\Delta m/m^-$。库存非负、有限且每节点至多支取一次；同一加入量不在入口预富集和沿程库存中重复计量。

机制对照M0关闭新增水但保留岩屑，M1加入同一有限水库存并计入静止混合，M2与M1的水岩库存完全相同，仅额外把累计新增水严格大于零后的下游动摩擦乘以0.6。M2的系数0.6是预先设定的假设映射，不是浓度率定律，也不依据是否到达口岸进行选择。三种对照的运程差见表\ref{tab:q4-inventory}；表中只比较同一剖面、同一事件定义下的正速到达结果，未达组不参与时间差计算。

\begin{table}[htbp]\centering\small
\caption{有限库存机制对照}\label{tab:q4-inventory}
\begin{tabular}{@{}c c c c@{}}
\toprule
$q_r$&$q_w$&$\Delta L_{\mathrm{M1-M0}}$/m&$\Delta L_{\mathrm{M2-M1}}$/m\\
\midrule
0.5&0.5&$-166.797$&$+4980.217$\\
0.5&1.0&$-608.791$&$+5422.211$\\
1.0&1.0&$-935.331$&$+6366.604$\\
\bottomrule
\end{tabular}
\end{table}

仅加入静止水会增加质量并按动量守恒降低速度，故M1相对M0缩短运程；M2的额外距离来自唯一改变的下游减阻映射。它不证明“水越多越远”，也不说明真实沿程侵蚀会按此发生。M0、M1和M2对照见图\ref{fig:q4-mechanism}。

\paperfig{figures/04c_合成补水机制.png}{0.80}{库存机制运程对照}{fig:q4-mechanism}

有限库存系列的数值检查包含23项合成测试，全部通过。当前最大质量相对残差为$2.076390489509933\times10^{-16}$，全程能量相对残差为$6.18040483706205\times10^{-15}$，段能量相对残差为$2.451294653897242\times10^{-12}$，节点动量和节点能量相对残差分别为$4.429633044287857\times10^{-16}$和$2.8349651483442284\times10^{-15}$。这些结果表明固定路径积分、库存扣减和解析停止规则与设定一致；物源可动性、真实流变或M2减阻映射仍需现场验证。

\subsection{综合判读与适用范围}
问题四的主要结论是：在给定体积、公共断面和参数下，一维基线于9.336769 km自然能量停止，未达口岸；D8的16个配置在17.5704376--44.7278067 min内到达口岸代理，但全部于26045.4946703 m分析域末端右删失，自然最大运程尚不能确定。两引擎差异来自停止规则、有效断面宽度、缓坡驱动与路由结构，不是同一观测量的两次测量。

D8质量账本说明离散转移守恒；地形时效、源区体积、屈服门槛、糙率和单流向假设是否正确，仍需现场检验。有限库存模型表明静止加水通常因质量稀释缩短运程，只有预先设定的减阻假设才显著延长，而该减阻缺少独立观测。未取得可核查原文和观测定义的公开报道量值不进入验证。因此，D8到达时间是条件路由结果，右删失末端是计算范围，M2是机制假设；升级为口岸风险预报仍需源区释放、沿程流深流速、可核实物源库存、现场流变和独立多事件观测。
